\pdfinfoomitdate=1
\pdftrailerid{}
\pdfsuppressptexinfo=15
\documentclass[11pt,letterpaper]{article}
\usepackage[T1]{fontenc}
\usepackage{lmodern}
\usepackage[margin=1in]{geometry}
\usepackage{amsmath,amssymb,amsthm,mathtools,booktabs,array}
\usepackage{microtype,xcolor,enumitem,fancyhdr}
\usepackage[colorlinks=true,linkcolor=black,citecolor=black,urlcolor=blue!55!black]{hyperref}
\hypersetup{pdftitle={Leading equivalents for self complementary tournament score sequences},pdfsubject={Report 116: A345470 and A351869},pdfauthor={Research report}}
\setlength{\headheight}{14pt}
\pagestyle{fancy}\fancyhf{}
\fancyhead[L]{\small Self-complementary score sequences}
\fancyhead[R]{\small Report 116}
\fancyfoot[C]{\thepage}
\renewcommand{\headrulewidth}{0.3pt}
\setlength{\parskip}{4pt}\setlength{\parindent}{0pt}
\setlist{topsep=4pt,itemsep=3pt,parsep=0pt}
\emergencystretch=2em\allowdisplaybreaks[1]
\newtheorem{theorem}{Theorem}[section]
\newtheorem{lemma}[theorem]{Lemma}
\newtheorem{proposition}[theorem]{Proposition}
\newtheorem{corollary}[theorem]{Corollary}
\theoremstyle{definition}\newtheorem{remark}[theorem]{Remark}
\newcommand{\Pp}{\mathbb P}\newcommand{\Ee}{\mathbb E}
\newcommand{\Z}{\mathbb Z}\newcommand{\R}{\mathbb R}
\newcommand{\ind}{\mathbf 1}\newcommand{\Vdw}{V_{\mathrm{DW}}}
\newcommand{\fY}{f_Y}
\title{\vspace{-1.2em}\textbf{Leading equivalents for self complementary\\tournament score sequences}}
\author{Report 116}\date{2 October 2026}
\begin{document}\maketitle\thispagestyle{empty}
\begin{abstract}
Let $C_n$ count self-complementary tournament score sequences of length $n$
(OEIS A345470), and let $D_n$ count the strong ones (A351869). We derive
\[
 C_n\sim A\,2^n n^{-3/4},\qquad
 D_n\sim e^{-\lambda}A\,2^n n^{-3/4},
\]
with one positive amplitude $A$ for both parities. The amplitude has an exact
harmonic-function characterization; no numerical value for it is established.
The limiting strong fraction satisfies
$0.7187529762<e^{-\lambda}<0.7187529792$, using Kolesnik's published bounds.
The new probabilistic step upgrades a survival weak limit to a fixed-velocity
endpoint equivalent by smoothing with a short final block. The proof treats
marginal regularity, the killing error, final slack and convolution tails.
It also gives a harmonic identity, a Lambert $W$ smooth inverse, sharp
vanishing-width ceiling brackets for integer thresholds, and an eventual
zero-or-one threshold gap, with both values occurring infinitely often.
No convergence rate or higher-order counting correction is proved.
\end{abstract}

\section{Result and conventions}\label{sec:result}
A tournament score sequence is the sorted list of vertex outdegrees.
Self-complementarity of a score sequence means
\begin{equation}\label{eq:complement}
 s_i+s_{n+1-i}=n-1\qquad(1\le i\le n).
\end{equation}
We count distinct sorted sequences, not labeled tournaments or tournament
isomorphism classes. A sequence is \emph{strong} when its tournament
realizations are strongly connected. Landau's criterion and its strict
version characterize these two classes by nonnegative and positive proper
partial-sum excesses, respectively; see \cite[Propositions 1 and 7]{stockmeyer}.

Let $C_n$ and $D_n$ denote the all and strong self-complementary counts.
Set $C_0=D_0=1$ for the empty terms in generating functions. The small values
are $C_1=D_1=1$ and $C_2=1,D_2=0$; empty-term conventions do not assert that an
empty tournament is strongly connected.

\begin{theorem}\label{thm:main}
There exists $A>0$ such that, through all positive integers $n$,
\begin{equation}\label{eq:main}
 C_n\sim A\,2^n n^{-3/4},\qquad
 D_n\sim rA\,2^n n^{-3/4},\qquad r=e^{-\lambda},
\end{equation}
where
\begin{align}
 \lambda&=\sum_{k\ge1}\frac{N_k}{k4^k},\label{eq:lambda}\\
 N_k&=\frac1{2k}\sum_{d\mid k}(-1)^{k+d}\varphi(k/d)\binom{2d}{d}.
 \label{eq:EGZ}
\end{align}
Here $\varphi$ is Euler's totient. In particular,
\begin{equation}\label{eq:ratio-bounds}
 0.7187529762<r<0.7187529792.
\end{equation}
For the variance-one geometric-minus-one walk defined in
Section~\ref{sec:DW}, with its Denisov--Wachtel normalization,
\begin{equation}\label{eq:amplitude}
 A=2^{1/4}\kappa\Vdw(1/\sqrt2,1/\sqrt2)\fY(0)>0.
\end{equation}
Thus the same leading constants hold on even and odd lengths.
\end{theorem}

Report 114 established the two $\Theta(2^n n^{-3/4})$ estimates, resolving the
order-of-growth questions in \cite[Section 6]{stockmeyer}. This separate report
strengthens that conclusion and recalls the encoding so that its proof can be
read independently. The underlying ordinary-score theorems, generating
functions and numerical bound on $\lambda$ are published results and are
attributed below. The application here concerns the self-complementary
counts. No claim of publication, peer review or priority is made.

\subsection{Proof outline and scope}
An exact geometric-gap encoding gives $2^{-n}C_n$ as a survival probability
with a final velocity constraint. Denisov--Wachtel Theorem 1 supplies survival
of order $n^{-1/4}$ and a conditional weak limit. A free last block supplies
one-dimensional endpoint mass of order $n^{-1/2}$. The substantive step is
to obtain its \emph{constant} while controlling the paths killed in that block.
Bounded continuity of the limiting velocity marginal and an endpoint-weighted
maximal estimate make this possible. A summable slack mixture then proves
the first equivalent. Stockmeyer's exact convolution and the ordinary-score
tail bound give the second. The proof uses no killed local-limit formula,
and specifically does not use equation (12) of \cite{dw}.

\section{Exact half-score and slack encoding}\label{sec:encoding}
Write $n=2m$ or $n=2m+1$ and
\[
 M=\left\lfloor\frac{n-1}{2}\right\rfloor
 =\begin{cases}m-1,&n=2m,\\m,&n=2m+1.\end{cases}
\]
For $m\ge1$, every self-complementary candidate is uniquely determined by
\begin{equation}\label{eq:half}
 0\le s_1\le\cdots\le s_m\le M.
\end{equation}
Reflect by \eqref{eq:complement}, inserting $s_{m+1}=m$ in the odd case.
The result is nondecreasing and its total is $\binom n2$.
For the full sequence put
\[
 A_j=\sum_{i=1}^j s_i-\binom j2\qquad(0\le j\le n).
\]
The last $j$ scores sum to $j(n-1)-\sum_{i\le j}s_i$, and therefore
\[
 A_{n-j}=\binom n2-j(n-1)+\sum_{i\le j}s_i-\binom{n-j}2=A_j.
\]
Consequently all Landau inequalities reduce exactly to
\begin{equation}\label{eq:Landau}
 A_j\ge0\quad(1\le j\le m);
 \qquad\text{strongness replaces $\ge0$ by $>0$.}
\end{equation}
This includes the middle inequality, and both middle inequalities for odd $n$.

Define the gaps and the \emph{last slack}
\begin{equation}\label{eq:gaps}
 G_1=s_1,\quad G_j=s_j-s_{j-1}\ (2\le j\le m),\quad
 G_{m+1}=M-s_m.
\end{equation}
This is a bijection between \eqref{eq:half} and nonnegative compositions
$\sum_{j=1}^{m+1}G_j=M$. Take the $G_j$ independent with
\begin{equation}\label{eq:geom}
 \Pp(G_j=g)=2^{-g-1},\qquad g\in\Z_{\ge0},
 \qquad X_j=G_j-1.
\end{equation}
Every such composition has probability $2^{-M-m-1}=2^{-n}$.
With initial velocity $V_0=1$ and area $A_0=0$, set
\[
 V_j=1+\sum_{i=1}^jX_i,\qquad A_j=\sum_{i=1}^jV_i.
\]
For $j\le m$, $V_j=s_j-j+1$, so $A_j$ is precisely the Landau excess.
The composition condition is equivalent to
\begin{equation}\label{eq:rn}
 V_{m+1}=r_n,\qquad
 r_n=\begin{cases}-1,&n=2m,\\0,&n=2m+1.\end{cases}
\end{equation}

For a general state $z=(a,v)$ we henceforth use
\begin{equation}\label{eq:process}
 V_j=v+\sum_{i=1}^jX_i,\qquad
 I_j=a+\sum_{i=1}^jV_i,\qquad
 \tau=\min\{j\ge1:I_j\le0\}.
\end{equation}
Because $A_j$ is integer, $A_j\ge0$ is equivalent to $I_j=A_j+1>0$.
Thus weak inequalities correspond to the start $(I_0,V_0)=(1,1)$.
Conditioning on the independent last slack $G_{m+1}=h$ gives the exact formula
\begin{equation}\label{eq:C-exact}
 2^{-n}C_n=\sum_{h\ge0}2^{-h-1}
 \Pp_{(1,1)}(\tau>m,\ V_m=r_n+1-h).
\end{equation}
The area at time $m+1$ is deliberately untested. For $n=2$, the valid weak
half sequence $(0)$ would fail a new test on its final unshifted area; for
$n=4$, the strong half sequence $(1,1)$ has final unshifted area zero and
would fail a strict extra test. The slack is not another score.

\section{The imported survival and weak-limit theorem}\label{sec:DW}
For \eqref{eq:geom},
\begin{equation}\label{eq:moments}
 \Ee X=0,\qquad \Ee X^2=\sigma^2=2,\qquad
 \Ee e^{\theta|X|}<\infty\quad(0<\theta<\log2).
\end{equation}
Its support $\{-1,0,1,\ldots\}$ has maximal lattice span one.
Apply Denisov--Wachtel \cite[Theorem 1, (7)--(8), pp.~169--170]{dw} to
$X/\sigma$. Its assumptions are mean zero, variance one and a finite
$(2+\delta)$-moment for some $\delta>0$, all verified by \eqref{eq:moments}.
The theorem's positive harmonic function $\Vdw$ is normalized as in that paper;
its universal survival constant is
\begin{equation}\label{eq:kappa}
 \kappa=\frac{3\Gamma(1/4)}{2^{3/4}\pi^{3/2}}.
\end{equation}
For a fixed point in
\[
 K_+=\{z\in(0,\infty)\times\R:
 \Pp_z(Z_j\in(0,\infty)^2,\tau>j)>0\text{ for some }j\ge0\},
\]
the theorem provides the survival equivalent and the conditional endpoint
weak limit. In this definition $Z_j$ denotes the area--velocity pair of the
variance-one process. Every positive-coordinate start belongs to $K_+$ at
$j=0$.

Divide \emph{both} coordinates of \eqref{eq:process} by $\sigma$. For any
fixed positive-coordinate integer state $z=(a,v)$ this gives
\begin{align}
 \Pp_z(\tau>n)&\sim c_z n^{-1/4},\qquad
 c_z=\kappa\Vdw(a/\sigma,v/\sigma)>0,\label{eq:survival}\\
 \mathcal L\left(\frac{I_n}{\sigma n^{3/2}},
       \frac{V_n}{\sigma\sqrt n}\ \middle|\ \tau>n\right)
 &\Longrightarrow\mu.\label{eq:weak-limit}
\end{align}
The measure $\mu$ is universal for these standardized walks. Theorem 1 has no
lattice condition; the span-one calculations below remain in the original
integer coordinates. In particular, the endpoint mass will contribute
$1/\sigma$, not a two-dimensional lattice-cell factor.

For precision, define $\Vdw$ using the paper's normalization
$\Vdw(z)=\lim_{n\to\infty}\Ee_z[h(Z_n);\tau>n]$, where $h$ is its Brownian
harmonic function \cite[equation (4)]{dw}. This fixes the scale of
\eqref{eq:amplitude}; it is not an arbitrary positive harmonic function.
Equivalently, the product $\kappa\Vdw(z)$ is fixed by \eqref{eq:survival}.

\section{Regularity of the limiting velocity density}\label{sec:marginal}
We must localize at velocity zero. Abstract weak convergence alone does not
supply this step. We verify the required regularity from the density displayed
in \cite[Theorem 1]{dw}. Write
\begin{align}
 \bar h(x,y)&=c_{\mathrm{DW}}\int_0^1\int_0^\infty
    w^{3/2}s^{-1/2}e^{-2w^2/s}
    q_{1-s}(x,-y;0,-w)\,dw\,ds,\label{eq:density}\\
 c_{\mathrm{DW}}&=\frac{2^{9/4}\sqrt\pi}{\sqrt3\Gamma(1/4)},\qquad
 q_t(x,y;u,v)=p_t(x,y;u,v)-p_t(x,y;u,-v),\nonumber
\end{align}
where the free Gaussian density is
\begin{equation}\label{eq:gaussian}
 p_t(x,y;u,v)=\frac{\sqrt3}{\pi t^2}
 \exp\left(-\frac{6(u-x-ty)^2}{t^3}
 +\frac{6(u-x-ty)(v-y)}{t^2}-\frac{2(v-y)^2}{t}\right).
\end{equation}
For $a\ge0$ define
\begin{equation}\label{eq:marginal}
 F_a(y)=\int_{x>a}\bar h(x,y)\,dx,\qquad \fY(y)=F_0(y).
\end{equation}

\begin{lemma}\label{lem:marginal}
The functions $F_a(y)$ are jointly continuous in $(a,y)\in[0,\infty)\times\R$
and bounded by a single finite constant. Moreover $0<\fY(0)<\infty$.
\end{lemma}
\begin{proof}
Completing the square in the starting area $x$ gives
\begin{equation}\label{eq:gaussian-integral}
 \int_\R p_t(x,-y;0,\pm w)\,dx
   =(2\pi t)^{-1/2}\exp\left(-\frac{(y\pm w)^2}{2t}\right).
\end{equation}
Consequently
\[
 \int_{x>a}|q_t(x,-y;0,-w)|\,dx\le2(2\pi t)^{-1/2}.
\]
Multiplying by the weight in \eqref{eq:density} and integrating over $w$
yields a constant times
\[
 s^{-1/2}(1-s)^{-1/2}\int_0^\infty w^{3/2}e^{-2w^2/s}\,dw
 =C s^{3/4}(1-s)^{-1/2}.
\]
This is integrable over $0<s<1$. Thus Fubini's theorem applies absolutely to
the integrals defining $F_a$, and they share a bound independent of $a,y$.
For fixed $s,w$, the integral over $x>a$ of either Gaussian in
\eqref{eq:gaussian-integral} is a continuous Gaussian-tail function of $a,y$.
The same integrable majorant and dominated convergence prove joint continuity,
including at $a=0$.

Finally, for $x,w,t>0$,
\[
 \frac{p_t(x,0;0,-w)}{p_t(x,0;0,w)}
 =\exp(12wx/t^2)>1.
\]
All outer weights in \eqref{eq:density} are positive, so
$\bar h(x,0)>0$ for $x>0$. Hence $F_a(0)>0$ for every finite $a$, and the
claimed assertion at zero follows.
\end{proof}
This argument selects continuous representatives of the marginal densities;
it does not rely on their almost-everywhere values at one chosen point.


\subsection{An explicit universal factor}
The factor $\fY(0)$ admits a convergent special-function expression independent
of the discrete harmonic function:
\begin{equation}\label{eq:fY-explicit}
 \fY(0)=\frac{5\Gamma(7/4)^2}{96\sqrt{2\pi}}
 {}_3F_2\left(\begin{matrix}1,7/4,9/4\\3/2,4\end{matrix};\frac34\right).
\end{equation}
Here ${}_pF_q$ has its usual series definition
$\sum_{j\ge0}(a_1)_j\cdots(a_p)_jz^j/
((b_1)_j\cdots(b_q)_j j!)$, where $(a)_j$ is the rising factorial.
The series in \eqref{eq:fY-explicit} is positive and convergent.

To derive the formula, integration of the two completed Gaussian squares
over $x>0$ gives, with $u=1-s$,
\[
 \int_0^\infty q_u(x,0;0,-w)\,dx
 =\frac{e^{-w^2/(2u)}}{\sqrt{2\pi u}}
       \operatorname{erf}\!\left(\sqrt{\frac3{2u}}\,w\right).
\]
For $P,Q>0$, the resulting $w$ integral is
\begin{equation}\label{eq:erf-integral}
 \int_0^\infty w^{3/2}e^{-Pw^2}\operatorname{erf}(Qw)\,dw
 =\frac{Q\Gamma(7/4)}{\sqrt\pi P^{7/4}}
 {}_2F_1\left(\frac12,\frac74;\frac32;-\frac{Q^2}{P}\right).
\end{equation}
For example, insert
$\operatorname{erf}(Qw)=2Qw\pi^{-1/2}\int_0^1e^{-Q^2w^2t^2}\,dt$
and first evaluate the nonnegative gamma integral in $w$.
Substitute $P=(4-3s)/(2s(1-s))$, $Q=\sqrt{3/(2(1-s))}$.
Pfaff's transformation changes the hypergeometric factor to
\[
 \left(\frac{4-3s}{4}\right)^{7/4}
 {}_2F_1\left(1,\frac74;\frac32;\frac{3s}{4}\right).
\]
After cancellation of powers, \eqref{eq:density} therefore gives
\[
 \fY(0)=\frac{\Gamma(7/4)}{\sqrt{2\pi}\Gamma(1/4)}
 \int_0^1s^{5/4}(1-s)^{3/4}
 {}_2F_1\left(1,\frac74;\frac32;\frac{3s}{4}\right)\,ds.
\]
Termwise integration is justified by positivity and convergence at $3/4$.
The beta integral contributes $B(9/4,7/4)$ and the additional hypergeometric
parameters $9/4$ and $4$. Finally
$\Gamma(9/4)=5\Gamma(1/4)/16$ and $\Gamma(4)=6$ give
\eqref{eq:fY-explicit}. This exact evaluation of the universal factor does
not evaluate $A$, which still contains $\Vdw(1/\sqrt2,1/\sqrt2)$.

\section{Free endpoint masses and a maximal estimate}\label{sec:bridge}
Put $T_j=X_1+\cdots+X_j$. The ordinary span-one lattice local central limit
theorem gives, uniformly over integers $d$,
\begin{equation}\label{eq:LCLT}
 \sqrt\ell\,\Pp(T_\ell=d)=
 \frac1{\sigma\sqrt{2\pi}}
 \exp\left(-\frac{d^2}{2\sigma^2\ell}\right)+o(1).
\end{equation}
In particular $\sup_d\Pp(T_\ell=d)\le C\ell^{-1/2}$ for $\ell\ge1$.
The hypotheses here are just the finite variance and span one already checked.
For this particular increment law the exact atom is also explicit:
\begin{equation}\label{eq:NB}
 \Pp(T_\ell=d)=\binom{2\ell+d-1}{\ell-1}2^{-2\ell-d}\quad(d\ge-\ell),
\end{equation}
and it is zero otherwise. The ratio of consecutive positive masses is
$(2\ell+d)/(2(\ell+d+1))$, so $d=-1$ is a mode and the maximum is
$\frac12\binom{2\ell-2}{\ell-1}4^{-(\ell-1)}=O(\ell^{-1/2})$.
Thus the uniform atom bound can alternatively be checked directly by Stirling.

\begin{lemma}\label{lem:max}
There is a constant $C$, depending only on the increment law, such that for
$\ell\ge4$, $H>0$, and an integer $d$ with $|d|\le H/2$,
\begin{equation}\label{eq:max}
 \Pp\left(T_\ell=d,\ \max_{j\le\ell}|T_j|\ge H\right)
 \le C\frac{\sqrt\ell}{H^2}.
\end{equation}
\end{lemma}
\begin{proof}
Split at $h=\lfloor\ell/2\rfloor$. For a violation in the first half, condition
on the first $h$ increments and bound the remaining independent endpoint atom
by $C/\sqrt{\ell-h}$. Kolmogorov's inequality bounds the first-half maximum
probability by $\sigma^2h/H^2$, giving $C\sqrt\ell/H^2$.

For a violation at a time $j>h$, the equality $T_\ell=d$ implies
$|T_\ell-T_j|\ge H-|d|\ge H/2$. The reversed last $\ell-h$ increments
therefore have a partial sum of absolute value at least $H/2$. Condition on
that suffix, use the first-half endpoint-atom bound $C/\sqrt h$, and then
Kolmogorov's inequality for its reversed partial sums. This contributes
another $C\sqrt\ell/H^2$. Their sum proves the estimate.
\end{proof}
The estimate concerns joint endpoint mass. Dividing by an endpoint probability
would introduce an unnecessary, possibly nonuniform bridge constant.

\section{Smoothing to one fixed terminal velocity}\label{sec:smoothing}
\begin{theorem}\label{thm:endpoint}
For every fixed positive-coordinate integer state $z=(a,v)$ and every fixed
integer $b$,
\begin{equation}\label{eq:endpoint}
 \Pp_z(\tau>n,V_n=b)\sim\frac{c_z\fY(0)}{\sigma}\,n^{-3/4}.
\end{equation}
Also, uniformly over integer $b$,
\begin{equation}\label{eq:uniform-endpoint}
 \Pp_z(\tau>n,V_n=b)\le C_z(n+1)^{-3/4}\qquad(n\ge1).
\end{equation}
\end{theorem}
\begin{proof}
For \eqref{eq:uniform-endpoint}, split at $\lfloor n/2\rfloor$, retain survival
only in the first block, and use the free endpoint-atom bound in the second.
The survival equivalent \eqref{eq:survival}, enlarged to a global upper bound,
then gives $C_z n^{-1/4}n^{-1/2}$. Enlarge $C_z$ for small times.

For the equivalent fix $0<\varepsilon<1/2$ and put
\[
 k=\lfloor(1-\varepsilon)n\rfloor,\quad \ell=n-k,
 \quad \alpha=1-\varepsilon.
\]
All limits in $n$ below keep $\varepsilon$ fixed.

\textbf{Upper bound.}
Dropping the area tests after $k$ gives
\begin{equation}\label{eq:upper-Markov}
 \Pp_z(\tau>n,V_n=b)
 \le\Ee_z[\ind_{\{\tau>k\}}\Pp(T_\ell=b-V_k)].
\end{equation}
The uniform error in \eqref{eq:LCLT} is $o(\ell^{-1/2})$; after averaging it is
multiplied by $\Pp_z(\tau>k)=O(k^{-1/4})$. Under the scaling
$V_k=\sigma\sqrt k\,y$, the Gaussian test functions converge uniformly on
$\R$ to the Gaussian with exponent $-\alpha y^2/(2\varepsilon)$.
Thus \eqref{eq:survival}--\eqref{eq:weak-limit} show that $n^{3/4}$ times the
right side of \eqref{eq:upper-Markov} converges to
\begin{equation}\label{eq:Ueps}
 U_\varepsilon=\frac{c_z\alpha^{-1/4}}{\sigma\sqrt{2\pi\varepsilon}}
 \int_\R e^{-\alpha y^2/(2\varepsilon)}\fY(y)\,dy.
\end{equation}
Writing $t=\varepsilon/\alpha$ and
$\phi_t(y)=(2\pi t)^{-1/2}e^{-y^2/(2t)}$, this is
\[
 U_\varepsilon=\frac{c_z}{\sigma}\alpha^{-3/4}
 \int_\R\phi_t(y)\fY(y)\,dy\longrightarrow\frac{c_z\fY(0)}{\sigma}
 \quad(\varepsilon\downarrow0)
\]
by Lemma~\ref{lem:marginal}. This proves the limsup bound.

\textbf{Lower bound before the killing error.}
Fix $\delta,K>0$ and at time $k$ retain only the states
\begin{equation}\label{eq:retained}
 I_k>\delta n^{3/2},\qquad |V_k|<K\sqrt n.
\end{equation}
First ignore killing in the last block. The same calculation gives the limit
\begin{equation}\label{eq:Leps}
 L_{\varepsilon,\delta,K}=
 \frac{c_z\alpha^{-1/4}}{\sigma\sqrt{2\pi\varepsilon}}
 \int_{|y|<K/(\sigma\sqrt\alpha)}
 e^{-\alpha y^2/(2\varepsilon)}
 F_{\delta/(\sigma\alpha^{3/2})}(y)\,dy.
\end{equation}
Indeed, in conditional scaled coordinates the retained set has area threshold
$\delta/(\sigma(k/n)^{3/2})$ and velocity cutoff $K/(\sigma\sqrt{k/n})$.
The three limiting boundary lines are $\mu$-null because $\mu$ has a density.
Weak convergence therefore applies to the bounded Gaussian times this
indicator. Moving boundary strips caused by floors have vanishing mass by
Portmanteau. No uniform weak-limit rate in the parameters is being used.

\textbf{Endpoint mass of paths killed in the last block.}
For a retained state $(a,v)$, the area after $j$ more steps is
\[
 a+jv+\sum_{i=1}^jT_i.
\]
For large $n$, $\ell\le2\varepsilon n$. If
$\varepsilon\le\delta/(8(K+1))$, then a failure of positivity implies
\begin{equation}\label{eq:H}
 \max_{i\le\ell}|T_i|\ge H,\qquad
 H=\frac{\delta}{4\varepsilon}\sqrt n.
\end{equation}
In fact $a-\ell|v|>\delta n^{3/2}-2\varepsilon Kn^{3/2}
\ge\delta n^{3/2}/2$, so a nonpositive area requires
$\max|T_i|\ge\delta n^{3/2}/(2\ell)\ge H$.
The endpoint condition is $T_\ell=b-v$ and eventually
$|b-v|\le(K+1)\sqrt n\le H/2$. Lemma~\ref{lem:max} therefore gives,
uniformly over all retained states,
\begin{equation}\label{eq:kill-error}
 \Pp_{(a,v)}(\tau\le\ell,V_\ell=b)
 \le C\varepsilon^{5/2}\delta^{-2}n^{-1/2}.
\end{equation}
The constant here does not grow as $\varepsilon\downarrow0$; the threshold
from which $n$ is sufficiently large may depend on $b,\delta,K,\varepsilon$.
Multiplying by $\Pp_z(\tau>k)$ shows that the error in $n^{3/4}$-scaled mass
is at most $C_z\alpha^{-1/4}\varepsilon^{5/2}\delta^{-2}$ in the limit.
Consequently
\begin{equation}\label{eq:liminf}
 \liminf_{n\to\infty}n^{3/4}\Pp_z(\tau>n,V_n=b)
 \ge L_{\varepsilon,\delta,K}
      -C_z\alpha^{-1/4}\varepsilon^{5/2}\delta^{-2}.
\end{equation}

\textbf{Taking the limits in order.}
In \eqref{eq:Leps} substitute $y=\sqrt{\varepsilon/\alpha}\,u$ to obtain
\[
 L_{\varepsilon,\delta,K}=\frac{c_z}{\sigma}\alpha^{-3/4}
 \int_{|u|<K/(\sigma\sqrt\varepsilon)}
 \frac{e^{-u^2/2}}{\sqrt{2\pi}}
 F_{\delta/(\sigma\alpha^{3/2})}
       (\sqrt{\varepsilon/\alpha}\,u)\,du.
\]
For fixed $\delta,K>0$, joint continuity and the common bound of
Lemma~\ref{lem:marginal} give
$L_{\varepsilon,\delta,K}\to(c_z/\sigma)F_{\delta/\sigma}(0)$ as
$\varepsilon\downarrow0$. The error in \eqref{eq:liminf} vanishes.
Finally let $\delta\downarrow0$, giving $F_{\delta/\sigma}(0)\to\fY(0)$.
This agrees with the upper bound and proves \eqref{eq:endpoint}.
The order is $n\to\infty$, then $\varepsilon\downarrow0$ at fixed $\delta,K$,
then $\delta\downarrow0$. Any one fixed $K>0$ suffices.
\end{proof}

\section{The common amplitude for all self-complementary counts}\label{sec:C}
Apply Theorem~\ref{thm:endpoint} to each fixed $h$ in \eqref{eq:C-exact}.
For both terminal values $r_n=-1,0$,
\[
 m^{3/4}\Pp_{(1,1)}(\tau>m,V_m=r_n+1-h)
 \longrightarrow\frac{c_{(1,1)}\fY(0)}{\sigma}.
\]
The uniform bound \eqref{eq:uniform-endpoint} gives a summable majorant after
multiplication by $2^{-h-1}$. Dominated convergence and
$\sum_{h\ge0}2^{-h-1}=1$ give
\begin{equation}\label{eq:C-asymptotic}
 2^{-n}C_n\sim\frac{c_{(1,1)}\fY(0)}{\sigma}\,m^{-3/4}.
\end{equation}
Since $m\sim n/2$ along either parity, its amplitude is
\[
 A=\frac{2^{3/4}c_{(1,1)}\fY(0)}{\sigma}
   =2^{1/4}\kappa\Vdw(1/\sqrt2,1/\sqrt2)\fY(0).
\]
The half-length conversion contributes $2^{3/4}$, the span-one endpoint mass
contributes $1/\sigma$, and the final-slack weights contribute total mass one.
There is no residual parity multiplier. This proves the first assertion of
Theorem~\ref{thm:main}.

\section{The strong equivalent from an exact convolution}\label{sec:D}
Let $S_k$ and $T_k$ count ordinary all and strong tournament score sequences,
with $S_0=1,T_0=0$. Write their ordinary generating functions as $S(z),T(z)$,
and use $C(z),D(z)$ for the self-complementary counts. Stockmeyer's
Theorems 11 and 13 \cite{stockmeyer} give
\begin{align}
 S(z)&=\frac1{1-T(z)},\label{eq:ordinary-gf}\\
 D_n&=C_n-\sum_{1\le k\le\lfloor n/2\rfloor}T_k C_{n-2k}\quad(n\ge1),
 \label{eq:Stock-conv}\\
 D(z)&=C(z)(1-T(z^2)),\qquad C(z)=D(z)S(z^2).
 \label{eq:SC-gf}
\end{align}
These are formal power-series identities, so no analytic assumption is needed
at this stage. The established ordinary-score asymptotic in \cite{kolesnik}
in particular yields
\begin{equation}\label{eq:ordinary-tail}
 0\le T_k/4^k\le S_k/4^k\le C(k+1)^{-5/2}.
\end{equation}
This ordinary-score input is not a new result of the present report.

\subsection{Justifying the convolution limit}
Put $c_n=C_n/2^n$ and $t_k=T_k/4^k$. We have
$c_n\sim A n^{-3/4}$ and hence $c_n\le B(n+1)^{-3/4}$ for all $n\ge0$.
For fixed $K$,
\[
 n^{3/4}\sum_{k=1}^K t_k c_{n-2k}
 \longrightarrow A\sum_{k=1}^K t_k.
\]
For $K<k\le n/4$, the bound on $c_{n-2k}$ gives
\begin{equation}\label{eq:conv-tail-one}
 n^{3/4}\sum_{K<k\le n/4}t_k c_{n-2k}
 \le C\sum_{k>K}t_k\longrightarrow0\quad(K\to\infty),
\end{equation}
uniformly in large $n$. The remaining terms satisfy
\begin{align}
 n^{3/4}\sum_{n/4<k\le n/2}t_k c_{n-2k}
 &\le C n^{3/4}n^{-5/2}
       \sum_{0\le j\le n/2}(j+1)^{-3/4}\nonumber\\
 &=O(n^{-3/2}).\label{eq:conv-tail-two}
\end{align}
Only one parity of $j=n-2k$ occurs, and extending the sum only enlarges it.
Thus
\begin{equation}\label{eq:conv-limit}
 n^{3/4}\sum_{1\le k\le n/2}t_k c_{n-2k}
 \longrightarrow A\sum_{k\ge1}t_k.
\end{equation}
Using \eqref{eq:Stock-conv},
\begin{equation}\label{eq:D-amplitude}
 D_n\sim A(1-T(1/4))\,2^n n^{-3/4}
       =\frac{A}{S(1/4)}\,2^n n^{-3/4}.
\end{equation}
The ordinary tail bound makes $S(1/4)$ finite. Taking increasing real limits
in \eqref{eq:ordinary-gf} gives $1-T(1/4)=1/S(1/4)>0$.

\subsection{Identifying the ratio and its numerical enclosure}
The ordinary recurrence and its generating-function form, recorded in
\cite{kolesnik,bdk}, are
\begin{equation}\label{eq:ordinary-recurrence}
 kS_k=\sum_{j=1}^kN_jS_{k-j},\qquad
 S(z)=\exp\left(\sum_{k\ge1}\frac{N_k}{k}z^k\right),
\end{equation}
with the exact divisor formula \eqref{eq:EGZ}. Its terms at $z=1/4$ are
$O(k^{-5/2})$: the $d=k$ term has this order after division by $k4^k$,
and all proper divisors satisfy $d\le k/2$, giving an exponentially smaller
bound. Thus the increasing real limit gives $S(1/4)=e^\lambda$ and proves
\eqref{eq:main} with $r=e^{-\lambda}$.

Kolesnik \cite[Section 3.1, equation (8)]{kolesnik} obtains the certified interval
\begin{equation}\label{eq:published-lambda}
 0.330237542\le\lambda\le0.330237546.
\end{equation}
Exponentiating with the reversed order yields \eqref{eq:ratio-bounds}.
The archived checker also independently reproduces a narrower enclosure by
computing $N_1,\ldots,N_{1669}$ exactly and combining their rational partial
sum with the tail estimate of \cite[Lemma 11, equation (17)]{kolesnik}.
Writing $K=1669$, that imported estimate is
\[
 \frac{1-2/K}{3\sqrt\pi K^{3/2}}
 \le \sum_{k>K}\frac{N_k}{k4^k}
 \le \frac1{3\sqrt\pi K^{3/2}}.
\]
Rational bounds for $\pi$, square roots and the exponential, with outward
rounding, certify in particular
\begin{align*}
 0.3302375420437401&\le\lambda\le0.3302375453489038,\\
 0.7187529767248666&\le r\le0.7187529791004629.
\end{align*}
The exact partial sum and full outward enclosures appear in
\path{checks/results.json}. This certificate uses a published infinite-tail
theorem; finite enumeration alone would not certify $\lambda$. We have
$1/2<r<1$ and $D_n/C_n\to r$.

These numerical enclosures concern $\lambda$ and $r$. Formula
\eqref{eq:amplitude} characterizes $A$ exactly, but contains a discrete harmonic
function that has not been evaluated or numerically enclosed here. The
limiting density is explicit; that does not by itself numerically determine
$A$.

\section{A direct strong amplitude and a harmonic identity}\label{sec:harmonic}
This section gives an independent expression for the strong amplitude.
It is optional for Theorem~\ref{thm:main}, whose proof is already complete.
It also checks that a mixture over unbounded starting states can be justified
without assuming a uniform version of \eqref{eq:survival}.

\begin{lemma}\label{lem:starts}
There is a constant $C$ such that, for all integers $g,m\ge1$,
\begin{align}
 \Pp_{(g,g)}(\tau>m)&\le C\sqrt g\,m^{-1/4},\label{eq:start-surv}\\
 \sup_{b\in\Z}\Pp_{(g,g)}(\tau>m,V_m=b)
   &\le C\sqrt g\,m^{-3/4}.\label{eq:start-atom}
\end{align}
In particular $c_{(g,g)}\le C\sqrt g$.
\end{lemma}
\begin{proof}
From $(1,1)$, use time $N=g^2$ and the normalized rectangle
\[
 2<I_N/(\sigma N^{3/2})<3,\qquad
 1<V_N/(\sigma\sqrt N)<2.
\]
It has positive $\mu$-mass. Indeed,
\[
 \frac{p_t(x,-y;0,-w)}{p_t(x,-y;0,w)}
   =\exp\left(\frac{4w(3x-yt)}{t^2}\right)>1
\]
for $0<t<1$ on this rectangle, so \eqref{eq:density} is positive there.
By \eqref{eq:survival}--\eqref{eq:weak-limit}, the probability of reaching it
while surviving is at least $c g^{-1/2}$ for all sufficiently large $g$.
Every state there dominates $(g,g)$ coordinatewise. For the finitely many
smaller $g$, force $G_1=g$ and $G_j=1$ for $2\le j\le g^2$.
Starting from $(1,1)$ the velocity then equals $g$ from time one onward and
the area stays positive; its time-$g^2$ state dominates $(g,g)$.
Each forced path has positive probability, so decreasing $c$ covers all $g$.

The common-increment coupling is monotone in both starting coordinates.
The Markov property therefore gives
\[
 \Pp_{(1,1)}(\tau>g^2+m)
 \ge c g^{-1/2}\Pp_{(g,g)}(\tau>m).
\]
The global upper bound following from \eqref{eq:survival} proves
\eqref{eq:start-surv}. Split an endpoint event at halfway and use the free
atom bound to obtain \eqref{eq:start-atom}, enlarging $C$ for $m=1$.
Taking $m\to\infty$ in the first bound gives the assertion about $c_{(g,g)}$.
\end{proof}

For strict Landau inequalities the first geometric value must be $g\ge1$.
At time one its unshifted state is $(A_1,V_1)=(g,g)$.
Decomposing this first value and the final slack yields, for $m\ge2$,
\begin{equation}\label{eq:D-direct}
 2^{-n}D_n=\sum_{g\ge1}\sum_{h\ge0}2^{-g-1}2^{-h-1}
 \Pp_{(g,g)}(\tau>m-1,V_{m-1}=r_n+1-h).
\end{equation}
Here the probabilities use the state after that first step as their new start.
Theorem~\ref{thm:endpoint} applies for each fixed $g,h$;
Lemma~\ref{lem:starts} and $\sum_{g\ge1}2^{-g-1}\sqrt g<\infty$ justify
summation. Hence the direct strong amplitude is
\begin{equation}\label{eq:B-direct}
 B=\frac{2^{3/4}\fY(0)}{\sigma}
     \sum_{g\ge1}2^{-g-1}c_{(g,g)}.
\end{equation}
Comparing with $B=rA$ gives the exact identity
\begin{equation}\label{eq:harmonic-identity}
 \sum_{g\ge1}2^{-g-1}\Vdw(g/\sqrt2,g/\sqrt2)
   =e^{-\lambda}\Vdw(1/\sqrt2,1/\sqrt2).
\end{equation}
The first-step weights are unconditional; their sum over $g\ge1$ is $1/2$.
Renormalizing them would incorrectly double the amplitude. The identity is a
consequence of two independently justified amplitude representations; whether
it has a shorter direct harmonic or combinatorial proof is a further question.

\section{Smooth inverses and integer thresholds}\label{sec:inverse}
All logarithms in this section are natural. Put $a=3/4$, $\gamma=\log2$,
and $A_C=A$, $A_D=rA$. For $E=C$ or $D$, define
\[
 F_E(t)=A_E e^{\gamma t}t^{-a}\quad(t>0),\qquad
 N_E(X)=\min\{n\ge1:E_n\ge X\}.
\]
The counts tend to infinity, so the minimum exists. Their equivalents imply
$E_{n+1}/E_n\to2$, hence eventual strict increase; a finite initial exception
such as $D_1>D_2$ has no effect for large $X$.

\begin{proposition}\label{prop:inverse}
For all sufficiently large $X$, the large real solution of $F_E(t)=X$ is
\begin{equation}\label{eq:Lambert}
 q_E(X)=-\frac a\gamma W_{-1}\left(
       -\frac\gamma a\left(\frac{A_E}{X}\right)^{1/a}\right),
\end{equation}
where $W_{-1}$ is the real branch with values at most $-1$. It satisfies
\begin{equation}\label{eq:inverse-expand}
 q_E(X)=\frac{\log X}{\gamma}
 +\frac a\gamma\log\log X
 -\frac a\gamma\log\gamma
 -\frac{\log A_E}{\gamma}+o(1).
\end{equation}
Moreover $q_E(E_n)=n+o(1)$. There exists $\eta_E(X)\ge0$ with $\eta_E(X)\to0$ such that, for all
sufficiently large $X$,
\begin{equation}\label{eq:ceil-brackets}
 \left\lceil q_E(X)-\eta_E(X)\right\rceil
 \le N_E(X)\le
 \left\lceil q_E(X)+\eta_E(X)\right\rceil.
\end{equation}
In particular $N_E(X)=q_E(X)+O(1)$, but a global claim
$N_E(X)=q_E(X)+o(1)$ is not justified and is false.
\end{proposition}
\begin{proof}
The function $F_E$ is strictly increasing for $t>a/\gamma$.
Rearranging its equation gives
$t e^{-\gamma t/a}=(A_E/X)^{1/a}$, from which \eqref{eq:Lambert} follows.
The branch $W_{-1}$ gives the large solution; $W_0$ would give a solution
approaching zero.

Let $L=\log X$ and set
\[
 u=\frac1\gamma(L+a\log L-a\log\gamma-\log A_E).
\]
Then $u\sim L/\gamma$ and
$\log F_E(u)=\log A_E+\gamma u-a\log u=L+o(1)$.
Since $(\log F_E)'(t)=\gamma-a/t$ is bounded away from zero at large $t$,
the mean value theorem gives $q_E(X)-u=o(1)$, proving
\eqref{eq:inverse-expand}. Similarly
$\log E_n-\log F_E(n)=o(1)$ proves $q_E(E_n)=n+o(1)$.

For the integer statement write $e_n=q_E(E_n)-n\to0$ and, for large $X$,
let $m=N_E(X)$. Eventual strict increase and removal of the finite initial
range give $E_{m-1}<X\le E_m$, hence
\[
 m-1+e_{m-1}<q_E(X)\le m+e_m.
\]
Take, for example, $\eta_E(X)=\max(|e_{m-1}|,|e_m|)$.
It tends to zero because $m\to\infty$, and
$q_E(X)-\eta_E(X)\le m<q_E(X)+1+\eta_E(X)$.
The two ceiling inequalities follow, including the case when a bound is an
integer. These brackets are qualitative and not an effective computable error
bound. To see why an $o(1)$ error for $N_E$ cannot hold for all real $X$, put
$X=F_E(j+1/2)$. Then $q_E(X)=j+1/2$, while its distance from any integer is
$1/2$. For large $j$, the brackets in fact force $N_E(X)=j+1$.
\end{proof}

In base-two form the smooth expansion is
\[
 q_E(X)=\log_2X+\frac34\log_2\log_2X-\log_2A_E+o(1).
\]
In particular
\begin{equation}\label{eq:smooth-gap}
 q_D(X)-q_C(X)\longrightarrow\frac\lambda{\log2}\in(0,1).
\end{equation}
The nonintegral limiting smooth displacement must not be confused with the
integer threshold gap.

\begin{corollary}\label{cor:gap}
For every sufficiently large real $X$,
\begin{equation}\label{eq:threshold-gap}
 N_C(X)\le N_D(X)\le N_C(X)+1.
\end{equation}
Both possible gaps, zero and one, occur for infinitely many $X$ tending to
infinity, even if thresholds are restricted to integers.
\end{corollary}
\begin{proof}
Always $D_n\le C_n$, and \eqref{eq:main} gives
$D_{n+1}/C_n\to2r>1$. For $m=N_C(X)$ large, therefore
$D_{m+1}>C_m\ge X$, so $N_D(X)\le m+1$; the other inequality is immediate.
At $X=D_n$ the eventual inequalities
$C_{n-1}<D_n\le C_n$ and strict growth of both counts give
$N_C(X)=N_D(X)=n$. At $X=C_n$, we have eventually
$D_n<C_n<D_{n+1}$, giving $N_C(X)=n$ and $N_D(X)=n+1$.
Both sequences of thresholds are integer and tend to infinity.
\end{proof}
Neither a rate for \eqref{eq:main} nor a higher-order coefficient expansion is
used in these arguments. Further terms of the smooth Lambert inverse alone
would not supply corresponding further terms for the actual count thresholds.

\section{Exact checks and reproducibility}\label{sec:checks}
The accompanying reader archive separates finite verification from the
analytic argument. Exact finite counts can check an encoding, coefficient
identity or decimal transformation; they cannot prove the survival theorem,
its limiting density, or the leading equivalents.

The package contains the editable TeX, rendered PDF, frozen reference
fixtures, executable arithmetic checks, mutation tests, a strict hash manifest,
and scripts for replay and deterministic rebuilding. Source-paper PDFs and
working research notes are not included. Primary references are linked below.
The finite-check coverage and ranges are specified in
\texttt{checks/reader\_validation.md}, and the actual results are preserved
in machine-readable files.

The reader workflow is:
\begin{enumerate}
\item Run \texttt{python3 replay.py}. It checks the exact declared file set,
then runs the arithmetic, semantic/schema mutation and manifest-attack suites
both normally and under \texttt{python3 -O}, comparing every result with the
archived output.
\item Run \texttt{python3 replay.py --build-pdf} to additionally rebuild the
PDF twice in independent temporary directories and compare both builds with
the archived PDF byte for byte. The recorded TeX toolchain is needed for exact
byte identity; another version may render the same mathematics differently.
\item Run \texttt{python3 make\_archive.py --output /tmp/report116.zip} from an
unchanged extracted package to recreate its deterministic ZIP bytes.
\end{enumerate}
The scripts use explicit checks that remain active with optimization. Replay
does not require network access and does not modify the sealed package.
The manifest rejects unexpected or missing files and directories, links,
special files, duplicate or unsafe paths, malformed entries and hash changes.
It is an integrity record, not a signature or a mathematical certificate.
Theorem~\ref{thm:main} rests on the proof and its stated published inputs.

\section{Further questions and precise limits}\label{sec:questions}
\begin{enumerate}
\item \textbf{Evaluate and certify the amplitude.}
Develop rigorous bounds for $\Vdw(1/\sqrt2,1/\sqrt2)$ in the specified
normalization and combine them with controlled quadrature for $\fY(0)$.
The exact expression \eqref{eq:amplitude} and the harmonic identity
\eqref{eq:harmonic-identity} provide consistency checks. A fit to moderate
counts would be evidence, not a certified value.
\item \textbf{Obtain quantitative errors and higher corrections.}
The smoothing proof takes successive limits and imports weak convergence
without a rate. A rate would require quantitative versions of the survival
and conditional limit inputs, together with corresponding marginal and
last-block estimates. Any correction analysis must track the two terminal
values, the half-length conversion and the slack mixture separately before
claiming that correction terms are parity-independent.
\item \textbf{Explain the harmonic identity directly.}
Find a potential-theoretic or combinatorial proof of
\eqref{eq:harmonic-identity} that does not compare two asymptotic amplitude
formulas. Such a proof could clarify its computational value or possible
extensions to other increment laws and symmetry constraints.
\item \textbf{Make inverse bounds effective.}
Certified bounds for the amplitude and counting error could turn
\eqref{eq:ceil-brackets} into usable threshold brackets. No explicit
rounding cutoff is supplied here.
\end{enumerate}
The established conclusions are the full leading equivalents, their common
parity constants, the exact limiting strong fraction, the stated harmonic
identity, and the qualitative inverse refinements. No numerical amplitude,
convergence rate, higher-order counting coefficient or effective asymptotic
threshold is asserted.

\begin{thebibliography}{9}
\bibitem{dw}
D. Denisov and V. Wachtel,
\emph{Exit times for integrated random walks},
Annales de l'Institut Henri Poincar\'e, Probabilit\'es et Statistiques
\textbf{51} (2015), 167--193.
Theorem 1, equations (7)--(8), and its displayed density on pp.~169--170.
\href{https://doi.org/10.1214/13-AIHP577}{doi:10.1214/13-AIHP577}.
\href{https://www.numdam.org/item/10.1214/13-AIHP577.pdf}{Open full text}.

\bibitem{stockmeyer}
P. K. Stockmeyer,
\emph{Counting Various Classes of Tournament Score Sequences},
Journal of Integer Sequences \textbf{26} (2023), Article 23.5.2.
Propositions 1 and 7, Theorems 11 and 13, and Section 6.
\href{https://cs.uwaterloo.ca/journals/JIS/VOL26/Stockmeyer/stock14.pdf}{Open full text}.

\bibitem{kolesnik}
B. Kolesnik,
\emph{The Asymptotic Number of Score Sequences},
Combinatorica \textbf{43} (2023), 827--844.
Ordinary-score asymptotics, generating functions, Lemma 11, equation (17),
and Section 3.1, equation (8), for the interval \eqref{eq:published-lambda}.
\href{https://doi.org/10.1007/s00493-023-00037-4}{doi:10.1007/s00493-023-00037-4}.

\bibitem{bdk}
M. Bassan, S. Donderwinkel and B. Kolesnik,
\emph{Tournament score sequences, Erd\H{o}s--Ginzburg--Ziv numbers,
and the L\'evy--Khintchine method},
Electronic Communications in Probability \textbf{31} (2026), paper 3.
Theorem 1.2, Proposition 1.3 and equations (1.2)--(1.5).
\href{https://doi.org/10.1214/26-ECP751}{doi:10.1214/26-ECP751}.

\bibitem{oeisC}
The OEIS Foundation, \emph{A345470: Self-complementary tournament score sequences}.
\href{https://oeis.org/A345470}{oeis.org/A345470}.

\bibitem{oeisD}
The OEIS Foundation, \emph{A351869: Strong self-complementary tournament score sequences}.
\href{https://oeis.org/A351869}{oeis.org/A351869}.
\end{thebibliography}
\end{document}
